% ======================================================================
\chapter{Exercícios práticos}
\label{cap:exercicios}
% ======================================================================

Faça os exercícios na ordem. Todos os comandos são digitados no \textbf{Git
Bash} (Windows) ou no \textbf{Terminal} (Linux e macOS).

Os materiais dos exercícios (\arq{jogo-da-memoria.zip}, \arq{main.py} e
\arq{modelo-actions.zip}) podem ser baixados pelos botões do topo da página
\url{https://andradasdev.github.io/github/} ou na pasta \arq{materials/} do
repositório \texttt{ronidomingues/github-training}. Os comandos abaixo supõem que eles foram
salvos na pasta \arq{Downloads}.

\begin{dica}[Pasta de treino]
Crie uma pasta só para os exercícios, por exemplo \arq{\textasciitilde/treino},
e entre nela com \comando{cd \textasciitilde/treino} antes de começar.
\end{dica}

% ----------------------------------------------------------------------
\section{Exercício 1 — Conferir as ferramentas}

\begin{bashcode}
git --version
gh --version
gh auth status
\end{bashcode}

\textbf{Como saber se deu certo:} as duas primeiras linhas mostram números de
versão, e a terceira diz \texttt{Logged in to github.com account
<seu-usuario>}. Se não estiver logado, volte ao capítulo~\ref{cap:git}.

% ----------------------------------------------------------------------
\section{Exercício 2 — O primeiro repositório}

\begin{enumerate}
  \item Crie o repositório no GitHub e já traga a cópia para o computador:
\begin{bashcode}
gh repo create meu-primeiro-repo --public --description "Meu primeiro repositório" --clone
cd meu-primeiro-repo
\end{bashcode}
  O Git avisa \texttt{You appear to have cloned an empty repository}. É
  normal: o repositório ainda não tem nenhum arquivo.

  \item Crie o \arq{README.md}, a \enquote{capa} do projeto no GitHub, escrita
        em Markdown:
\begin{bashcode}
cat > README.md << 'EOF'
# Meu Primeiro Projeto

Este é meu primeiro repositório no GitHub!

## O que este projeto faz?

- Aprender Git e GitHub
- Praticar comandos
EOF
\end{bashcode}

  \item Registre e envie:
\begin{bashcode}
git add README.md
git commit -m "docs: cria o README"
git push -u origin main
\end{bashcode}
\end{enumerate}

\begin{conceito}[O que o -u faz]
O \texttt{-u} liga a sua branch \texttt{main} à branch \texttt{main} do
GitHub (\texttt{origin}). Isso só precisa ser feito no primeiro push: depois,
basta \comando{git push} e \comando{git pull}.
\end{conceito}

\textbf{Como saber se deu certo:} a página
\texttt{github.com/<seu-usuario>/meu-primeiro-repo} mostra o README
formatado.

% ----------------------------------------------------------------------
\section{Exercício 3 — Uma licença criada pelo site}

Este exercício mostra o caminho inverso: uma mudança feita no GitHub que
precisa chegar ao seu computador.

\begin{enumerate}
  \item Na página do repositório, clique em \menu{Add file} \(\rightarrow\)
        \menu{Create new file}.
  \item Digite \texttt{LICENSE} como nome. Aparece o botão
        \menu{Choose a license template}: escolha \menu{MIT License},
        confira seu nome e clique em \menu{Review and submit}.
  \item Clique em \menu{Commit changes}.
  \item No terminal, traga a novidade:
\begin{bashcode}
git pull
ls
\end{bashcode}
\end{enumerate}

\textbf{Como saber se deu certo:} o \comando{ls} mostra o arquivo
\arq{LICENSE}, e o \comando{git log --oneline} tem dois commits.

% ----------------------------------------------------------------------
\section{Exercício 4 — Resolver um conflito}

\begin{enumerate}
  \item \textbf{No site:} abra o \arq{README.md} do
        \texttt{meu-primeiro-repo}, clique no lápis (\menu{Edit}), troque a
        primeira linha por \texttt{\# Meu Primeiro Projeto (editado no site)}
        e clique em \menu{Commit changes}.
  \item \textbf{No computador, sem dar pull:} abra o \arq{README.md} num
        editor, troque a \emph{mesma} primeira linha por
        \texttt{\# Meu Primeiro Projeto (editado no computador)}, salve e
        registre:
\begin{bashcode}
git commit -am "docs: edita o título no computador"
git push
\end{bashcode}
  O push é recusado, porque o GitHub tem um commit que você não tem:
\begin{saida}
 ! [rejected]        main -> main (fetch first)
error: failed to push some refs to 'https://github.com/...'
\end{saida}

  \item Traga a alteração do site:
\begin{bashcode}
git pull
\end{bashcode}
\begin{saida}
CONFLICT (content): Merge conflict in README.md
Automatic merge failed; fix conflicts and then commit the result.
\end{saida}

  \item Abra o \arq{README.md}. O Git colocou as duas versões dentro dele:
\begin{saida}
<<<<<<< HEAD
# Meu Primeiro Projeto (editado no computador)
=======
# Meu Primeiro Projeto (editado no site)
>>>>>>> 15261e4...
\end{saida}
  Entre \texttt{<<<<<<<} e \texttt{=======} está a sua versão; entre
  \texttt{=======} e \texttt{>>>>>>>}, a que veio do GitHub. Deixe só a
  linha que deve ficar (ou escreva uma terceira), apague os três marcadores e
  salve.

  \item Registre a resolução e envie:
\begin{bashcode}
git add README.md
git commit -m "docs: resolve o conflito do título"
git push
\end{bashcode}
\end{enumerate}

\textbf{Como saber se deu certo:} \comando{git log --oneline --graph} mostra
as duas linhas de trabalho se juntando no commit de resolução.

\begin{aviso}
Se o \comando{git pull} terminar em \texttt{Need to specify how to reconcile
divergent branches}, rode \comando{git config --global pull.rebase false} e
repita o pull.
\end{aviso}

% ----------------------------------------------------------------------
\section{Exercício 5 — Um site no GitHub Pages}

\begin{enumerate}
  \item Crie o repositório do site:
\begin{bashcode}
cd ~/treino
gh repo create meu-site --public --description "Meu primeiro site" --clone
cd meu-site
\end{bashcode}

  \item Descompacte o \arq{jogo-da-memoria.zip} dentro de \arq{Downloads}
        (no Windows: botão direito \(\rightarrow\) \menu{Extrair tudo}). Ele
        cria a pasta \arq{jogo-da-memoria/}. Copie o \textbf{conteúdo} dela,
        e não a pasta, para dentro do \arq{meu-site}:
\begin{bashcode}
cp -r ~/Downloads/jogo-da-memoria/. .
ls
\end{bashcode}
  O \comando{ls} precisa mostrar o \arq{index.html} aqui, na raiz.

  \item Registre e envie:
\begin{bashcode}
git add .
git commit -m "feat: adiciona o jogo da memória"
git push -u origin main
\end{bashcode}

  \item No GitHub: \menu{Settings} \(\rightarrow\) \menu{Pages}. Em
        \menu{Source}, escolha \menu{Deploy from a branch}; em
        \menu{Branch}, \texttt{main} e \texttt{/(root)}; clique em
        \menu{Save}.
\end{enumerate}

\textbf{Como saber se deu certo:} depois de um ou dois minutos, o jogo abre
em \texttt{https://<seu-usuario>.github.io/meu-site/}.

% ----------------------------------------------------------------------
\section{Exercício 6 — Automação com GitHub Actions}

\begin{enumerate}
  \item Crie o repositório e copie o \arq{main.py} para dentro dele:
\begin{bashcode}
cd ~/treino
gh repo create meu-script-python --public --description "Script Python com GitHub Actions" --clone
cd meu-script-python
cp ~/Downloads/main.py .
\end{bashcode}

  \item Crie a pasta dos workflows e o arquivo:
\begin{bashcode}
mkdir -p .github/workflows
\end{bashcode}
  Abra um arquivo novo \arq{.github/workflows/python.yml} no editor e cole o
  Exemplo 2 do capítulo~\ref{cap:pages}.

  \item Envie:
\begin{bashcode}
git add .
git commit -m "ci: roda o script Python a cada push"
git push -u origin main
\end{bashcode}

  \item Acompanhe na aba \menu{Actions} do repositório, ou com
        \comando{gh run watch}. Quando terminar, traga o arquivo gerado:
\begin{bashcode}
git pull
cat resultado.txt
\end{bashcode}
\end{enumerate}

\textbf{Como saber se deu certo:} o \arq{resultado.txt} mostra o horário da
execução, e o último commit do repositório é do \texttt{github-actions[bot]}.

\begin{dica}[Outro modelo]
O \arq{modelo-actions.zip}, em \arq{materials/}, traz o mesmo fluxo para um
programa em Fortran: compila, executa e grava o resultado.
\end{dica}

% ----------------------------------------------------------------------
\section{Exercício 7 — Fork, Pull Request e lista de presença}
\label{sec:ex_presenca}

A lista de presença deste guia (Anexo~\ref{an:presenca}) é feita por Pull
Request. É o fluxo real de contribuir com o projeto de outra pessoa.

\begin{enumerate}
  \item Crie o seu fork e clone:
\begin{bashcode}
cd ~/treino
gh repo fork ronidomingues/github-training --clone
cd github-training
\end{bashcode}

  \item Crie uma branch para a sua contribuição:
\begin{bashcode}
git switch -c presenca/seu-nome
\end{bashcode}

  \item Crie um arquivo com o seu nome na pasta \arq{presences/}. Use só
        letras, números, espaços, pontos e hífens no nome. Se quiser, escreva
        uma avaliação da capacitação dentro dele.
\begin{bashcode}
echo "Gostei muito da parte de conflitos!" > "presences/Seu Nome.txt"
\end{bashcode}

  \item Registre e envie para o \emph{seu} fork:
\begin{bashcode}
git add presences
git commit -m "chore(presences): registra a presença de Seu Nome"
git push -u origin presenca/seu-nome
\end{bashcode}

  \item Abra o Pull Request. No GitHub, na página do seu fork, clique em
        \menu{Compare \& pull request}, confira que a base é
        \texttt{ronidomingues/github-training} e clique em
        \menu{Create pull request}. Pelo terminal:
\begin{bashcode}
gh pr create --repo ronidomingues/github-training --fill
\end{bashcode}
\end{enumerate}

\textbf{Como saber se deu certo:} o PR aparece na aba \menu{Pull requests} do
repositório original. Depois que ele for aceito, o seu nome entra na lista
de presença da próxima versão do guia.

% ----------------------------------------------------------------------
\section{Exercício 8 — Mensagens de commit}

Reescreva cada mensagem no formato Conventional Commits
(capítulo~\ref{cap:padroes}):

\begin{enumerate}
  \item \enquote{arrumei o erro que dava quando a senha tava vazia}
  \item \enquote{coloquei explicação de como instalar no readme}
  \item \enquote{agora o relatório sai em pdf também}
  \item \enquote{atualizei o workflow pra rodar os testes}
\end{enumerate}

\begin{dica}[Respostas possíveis]
\texttt{fix(login): aceita senha vazia sem erro} ·
\texttt{docs(readme): explica a instalação} ·
\texttt{feat(relatorio): exporta em PDF} ·
\texttt{ci: roda os testes no workflow}
\end{dica}

% ----------------------------------------------------------------------
\section{Exercício 9 — Um commit assinado}

Siga o capítulo~\ref{cap:assinatura} com GPG ou SSH e faça um commit no
\texttt{meu-primeiro-repo}:

\begin{bashcode}
git commit --allow-empty -m "chore: testa a assinatura"
git push
git log --format='%h %G? %s' -1
\end{bashcode}

\textbf{Como saber se deu certo:} o terminal mostra a letra \texttt{G}, e a
aba \menu{Commits} do GitHub mostra o selo \textbf{Verified}.

% ----------------------------------------------------------------------
\section{Checklist}

\begin{itemize}
  \item[$\square$] Git e GitHub CLI instalados e logados
  \item[$\square$] Repositório criado, com README e LICENSE
  \item[$\square$] Um conflito resolvido
  \item[$\square$] Site publicado no GitHub Pages
  \item[$\square$] Workflow rodando no GitHub Actions
  \item[$\square$] Pull Request de presença aberto
  \item[$\square$] Commit assinado com selo \textbf{Verified}
\end{itemize}
